\section{Multivariable Polynomial Rings (Non-Examinable)}

A multivariate polynomial ring is not necessarily a PID. But how well can we understand these?

Let \(k\) be a field, and \(k\brac{x_1, \ldots, x_n}\) be a polynomial ring of \(n\) variables over \(k\). Given polynomials \(f_1, \ldots, f_s\), and another polynomial \(f\), (how) can we determine if
\[
    f \in \para{f_1, \ldots, f_s}?
\]

In the easy case (one-variable), the polynomial ring \(k\brac{x}\) is a PID. Given \(f_1, \ldots, f_s\), then
\[
    I = \para{f_1, \ldots, f_n} = \para{g}
\]
where
\[
    g = \gcd\para{f_1, \ldots, f_n}.
\]

But computing \(g\) is a totally algorithmic procedure (Euclidean Algorithm), and \(f \in \para{g}\) if and only if \(g \divides f\) -- and this is also completely algorithmic!

However, when we consider the multivariate case, even with the simple example where the ring is \(k\brac{x, y}\), and the ideal is
\[
    I = \para{x^2 + y^2 + 1, x - y},
\]
this question is non-trivial, since the ideals are not principle, and there is no Euclidean algorithm.

Surprisingly, the multivariate case \emph{does} have a solution!

We can define orderings for monomials (\emph{monomial orderings}), and for example, in two variables, we may declare
\[
    x^a y^b \succ x^c y^d
\]
if \(a > c\), or \(a = c\) and \(b = d\). This is called the \emph{lexicographic ordering}, or \emph{alphabetic ordering}. For example,
\[
    x^2 \succ x^2 y \succ x y^5 \succ y^{100}.
\]

This is not symmetric at all in \(x\) and \(y\), but this allows us to define the \emph{leading term} of every polynomial in \(k\brac{x, y}\), just like the single-variable case where we have a highest-power term. For example, we have
\[
    \text{LT}\para{3x^2 y + y^3 + 2} = 3x^2 y
\]
where we use \(\text{LT}\) to denote the leading term.

But once we made this choice, we have a \emph{multivariate version of the division algorithm}. Given \(f_1, \ldots, f_s\), and \(f\) in \(k\brac{x_1, \ldots, x_n}\), we want to come up with a good notion of division by the \emph{set} of polynomials. We want to write
\[
    f = a_1 f_1 + \cdots + a_s f_s + r
\]
where no term of \(r\) is divisible by any \(\text{LT} \para{f_i}\). (This reduces to the usual polynomial division in one variable and one polynomial.) Such a na\"{i}ve choice of the ordering, surprisingly, does allow this division to be possible.

This seems to completely solve the problem -- but not yet. If \(r = 0\), then \(f\) definitely sits in the ideal generated by the \(f_i\)s. However, if \(r \neq 0\), it is \emph{still possible} for \(f\) to sit in the ideal. The problem here is, when we are doing such divisions, all we are controlling is the leading terms, but leading terms might \emph{cancel} between different \(f_i\)s.

The key term here is Gr\"{o}bner Bases introduced by Bruno Buchberger (1960s) and Heisuke Hironaka (1960s). A generating set \(\brce{g_1, \ldots, g_t}\) for an ideal \(I\) is a \emph{Gr\"{o}bner basis} if every polynomial in \(I\) has leading terms divisible by some \(\text{LT}\para{g_i}\).

All we now have to show is, that there is some Gr\"{o}bner basis for any ideal, and in particular, we want to have some algorithm to find the Gr\"{o}bner basis for any ideal. Given \(f, g\), we define
\[
    S\para{f, g} = \frac{\lcm \para{\text{LT} \para{f}, \text{LT} \para{g}}}{\text{LT} \para{f}} \cdot f - \frac{\lcm \para{\text{LT} \para{f}, \text{LT} \para{g}}}{\text{LT} \para{g}} \cdot g.
\]

\emph{Buchberger's Algorithm} is as follows. We start with a generating set \(G = \brce{f_1, \ldots, f_s}\), and we repeat the following.

\begin{enumerate}
    \item For each \(f, g \in G\), we compute \(S\para{f, g}\).
    \item We divide \(S\para{f, g}\) by the set \(G\) (as defined above). If the remainder is non-zero, add that remainder to \(G\) (and we increase the size of \(G\)). We continue until all possible pairs are considered (including those within the newly-added ones).
\end{enumerate}

This algorithm terminates because of the Hilbert Basis Theorem.

This solves the ideal membership question perfectly! There is a theorem which states that, upon polynomial division by a Gr\"{o}bner Basis of some ideal, zero remainder is exactly the condition for the polynomial to be within the ideal. Formally, let \(I \idealsub k\brac{x_1, \ldots , x_n}\), and \(G\) be a Gr\"{o}bner Basis, then \(f \in I\), if and only if \(f\) gives remainder zero when divided by \(G\).